\documentclass[a4paper,10pt,landscape]{article}
\usepackage{multicol}
\usepackage{amsmath}
\usepackage{siunitx}
\usepackage{xcolor}
\usepackage{tcolorbox}
\usepackage[margin=0.7cm]{geometry}
\usepackage{graphicx}
\setlength{\columnsep}{1cm}
\setlength{\parindent}{0pt}

% ----- Couleurs -----
\definecolor{mechblue}{RGB}{0,90,180}
\definecolor{elecred}{RGB}{200,40,40}
\definecolor{graybg}{RGB}{245,245,245}
\definecolor{blockbg}{RGB}{255,255,255}

% ----- Style de boîte -----
\tcbset{
  cheatbox/.style={
    enhanced,
    sharp corners,
    colback=blockbg,
    colframe=black!40,
    boxrule=0.5pt,
    left=4pt, right=4pt, top=2pt, bottom=2pt,
    fonttitle=\bfseries,
    coltitle=black,
    before skip=6pt, after skip=6pt
  }
}

\begin{document}

\begin{center}
{\LARGE \textbf{Fiche Synthèse -- Énergie potentielle et Puissance}}\\[4pt]
\end{center}

\begin{multicols}{2}

% ----------------------------------------
% COLONNE 1 : MÉCANIQUE
% ----------------------------------------

\begin{tcolorbox}[cheatbox,title={\color{mechblue}Énergie potentielle gravitationnelle}]
\[
E = m g h
\]
$m$ : masse (\si{kg}) \quad
$g$ : gravité (\si{m/s^2}) \quad
$h$ : hauteur (\si{m}) \quad
$E$ : énergie (\si{J})

\textbf{Interprétation :}  
Chaque kilogramme placé à une hauteur $h$ possède une énergie potentielle $E=mgh$ qui peut être libérée en descendant.

\textbf{Exemple :}  
$1\,\si{m^3}$ d’eau ($1000\,\si{kg}$) à $100\,\si{m}$ :  
\[
E = 1000 \times 9.81 \times 100 = 981{,}000\,\si{J}
\]
→ environ 0.27 kWh stockés dans le lac.
\end{tcolorbox}


\begin{tcolorbox}[cheatbox,title={\color{mechblue}Puissance hydraulique}]
\[
P = \rho g h Q = \frac{E}{t}
\]
$P$ : puissance (\si{W}) \quad
$\rho$ : densité (\si{kg/m^3}) \quad
$g$ : gravité (\si{m/s^2}) \quad
$h$ : hauteur (\si{m}) \\
$Q$ : débit volumique (\si{m^3/s}) \quad
$t$ : temps (\si{s})

\textbf{Exemple :}  
Un barrage laisse passer $Q=1\,\si{m^3/s}$ d’eau depuis $100\,\si{m}$ :  
\[
P = 1000 \times 9.81 \times 100 \times 1 = 981{,}000\,\si{W} = 981\,\si{kW}
\]
→ la puissance mécanique transformée en électricité par la turbine.
\end{tcolorbox}


% ----------------------------------------
% COLONNE 2 : ÉLECTRICITÉ
% ----------------------------------------

\begin{tcolorbox}[cheatbox,title={\color{elecred}Énergie potentielle électrique (batterie)}]
\[
E = U Q = U I t
\]
$E$ : énergie (\si{J}) \quad
$U$ : tension (\si{V}) \quad
$Q$ : charge (\si{C}) \quad
$I$ : courant (\si{A}) \quad
$t$ : temps (\si{s})

\textbf{Interprétation :}  
Chaque coulomb (paquet d’électrons) placé à une tension $U$ possède une énergie potentielle $UQ$.  
Une batterie élève les charges “en hauteur électrique”, comme un barrage élève l’eau.

\textbf{Exemple :}  
Batterie $12\,\si{V}$, $10\,\si{Ah}$ :  
\[
10\,\si{Ah} = 10 \times 3600\,\si{C} = 36{,}000\,\si{C}
\]
\[
E = 12 \times 36{,}000 = 432{,}000\,\si{J} = 120\,\si{Wh}
\]
→ équivalent à soulever $60\,\si{kg}$ à 730 m.
\end{tcolorbox}


\begin{tcolorbox}[cheatbox,title={\color{elecred}Puissance électrique}]
\[
P = U I = \frac{E}{t}
\]
$P$ : puissance (\si{W}) \quad
$U$ : tension (\si{V}) \quad
$I$ : courant (\si{A}) \quad
$t$ : temps (\si{s})

\textbf{Exemple :}  
Une batterie $12\,\si{V}$ fournissant $10\,\si{A}$ :  
\[
P = 12 \times 10 = 120\,\si{W}
\]
\end{tcolorbox}

\end{multicols}


% ----------------------------------------
% ANALOGIE HYDRAULIQUE EN PLEINE LARGEUR
% ----------------------------------------

\begin{tcolorbox}[cheatbox, title={\color{elecred}Analogie mécanique – électrique}, width=0.97\textwidth]
\begin{center}
\begin{tabular}{lll}
\textbf{Mécanique (gravité)} & \textbf{Électricité} & \textbf{Sens physique}\\
\hline
Hauteur ($h$) & Tension ($U$) & Différence de potentiel\\
Masse ($m$) & Charge ($Q$) & Quantité stockée\\
Densité $\rho g$ & Capacité $C$ & “Rigidité” du stockage\\
Débit ($Q/t$) & Courant ($I$) & Flux de transfert\\
Force ($F = m g$) & Champ électrique ($E$) & Force motrice unitaire\\
\end{tabular}
\end{center}

\vspace{6pt}
\[
E = m g h \quad\Longleftrightarrow\quad E = U Q = U C V \\, 
P = F v = \frac{E}{t} \quad\Longleftrightarrow\quad P = U I = U \frac{Q}{t}
\]

\textbf{Idée clé :}  
Un barrage accumule de la mass à un potentiel gravitationnel de x mètres ;  
une batterie accumule des charges à un potentiel électrique de x Volts.  
Dans les deux cas, la puissance dépend du \textit{flux} (débit ou courant).
\end{tcolorbox}


% ----------------------------------------
% SECTION : SYNTHÈSE VISUELLE
% ----------------------------------------
\begin{center}
\begin{tcolorbox}[cheatbox, colback=graybg, width=0.97\textwidth]
\centering
\textbf{Cohérence dimensionnelle (analyse des unités du SI)}
\[
\underbrace{U I t}_{(\si{\volt\ampere\second})} =
\underbrace{E}_{(\si{\joule})} =
\underbrace{F d}_{(\si{\newton\metre})} =
\underbrace{m g h}_{(\si{\kilo\gram\metre\squared\per\second\squared})}
\]

\end{tcolorbox}
\end{center}

\end{document}
